% Presentación (Beamer) - Proyecto final de Matemáticas Discretas (MAT 4021), Grupo 7
% IMPORTANTE: reemplazar USUARIO/REPOSITORIO por el enlace real del repositorio (2 lugares) y recompilar.
\documentclass[aspectratio=169,14pt]{beamer}
\usepackage[utf8]{inputenc}
\usepackage[T1]{fontenc}
\usepackage[spanish]{babel}
\usepackage{amsmath,amssymb}
\usepackage{booktabs,colortbl,xcolor,tikz}
\usepackage{hyperref}
\usetikzlibrary{arrows.meta,positioning}

\definecolor{tinta}{HTML}{1B2A38}
\definecolor{acento}{HTML}{D9482B}
\definecolor{azul}{HTML}{2F6DB5}
\definecolor{papel}{HTML}{F6F4EF}
\definecolor{gris}{HTML}{5B6B7A}
\setbeamercolor{background canvas}{bg=papel}
\setbeamercolor{normal text}{fg=tinta}
\setbeamercolor{frametitle}{bg=tinta,fg=white}
\setbeamercolor{structure}{fg=azul}
\setbeamertemplate{navigation symbols}{}
\setbeamertemplate{frametitle}{\nointerlineskip
  \begin{beamercolorbox}[wd=\paperwidth,ht=3.4ex,dp=1.4ex,leftskip=1.2em]{frametitle}
  \usebeamerfont{frametitle}\insertframetitle\end{beamercolorbox}}
\setbeamerfont{frametitle}{series=\bfseries,size=\normalsize}
\setbeamertemplate{footline}{\hfill{\scriptsize\color{gris}\insertframenumber\,/\,\inserttotalframenumber}\hspace{1.2em}\vspace{3pt}}
\hypersetup{colorlinks=true,urlcolor=azul,linkcolor=azul}
\newcommand{\pendiente}[1]{\textbf{\textcolor{red}{[#1]}}}
\newcommand{\unidad}[1]{{\footnotesize\color{gris}#1}\par\vspace{2mm}}
\newcommand{\tarjeta}[2]{\begin{minipage}[t]{0.3\textwidth}\centering{\fontsize{28}{30}\selectfont\bfseries\color{acento}#1}\\[2mm]{\small #2}\end{minipage}}
\newcommand{\nodo}[1]{\tikz[baseline=-0.6ex]\node[draw=azul,rounded corners=2pt,fill=white,font=\scriptsize,inner sep=3pt,align=center,minimum width=1.25cm]{#1};}

\title{Asistente de ruta académica}
\subtitle{Modelación del plan de estudios con grafos, álgebra booleana y centralidad}
\author{Grupo 7: Juan Ayala Díaz \and Bradley De La Hoz Guerrero \and Álvaro León Curvelo \and Alejandro Márquez Rincón}
\institute{Matemáticas Discretas (MAT 4021) · Universidad del Norte}
\date{Noviembre de 2026}

\begin{document}

\begin{frame}[plain]
\titlepage
\vfill\hfill{\small Repositorio: \url{https://github.com/USUARIO/REPOSITORIO} \pendiente{enlace real}}
\end{frame}

\begin{frame}{El problema}
\unidad{Muchos estudiantes no saben qué se les bloquea}\vspace{2mm}
\tarjeta{65 \%}{no sabe con exactitud cuántos créditos ha aprobado}\hfill
\tarjeta{42 \%}{no sabe con seguridad qué materias se le bloquean si pierde o retira una}\hfill
\tarjeta{37 \%}{ha dudado entre retirar o perder sin saber las consecuencias}
\vfill
{\footnotesize\color{gris}Encuesta propia, anónima y voluntaria: 43 estudiantes de 17 programas (muestra por conveniencia; 84 \% cursa 2.º o 3.er semestre).}
\end{frame}

\begin{frame}{El plan de estudios es un grafo dirigido}
\unidad{Unidades 3 y 4 · grafo acíclico (DAG) de prerrequisitos}
\tarjeta{62}{vértices: materias, cada una con sus créditos}\hfill
\tarjeta{40}{aristas: 39 prerrequisitos y 1 correquisito}\hfill
\tarjeta{155}{créditos en 10 semestres (Ing. de Sistemas)}
\vspace{3mm}

\noindent {\small Cadena más larga de prerrequisitos (9 materias):}\\[1mm]
\resizebox{\textwidth}{!}{\begin{tikzpicture}[>=Stealth,every node/.style={draw=azul,rounded corners=2pt,fill=white,font=\footnotesize,align=center,minimum width=2.5cm,minimum height=0.8cm}]
\node (a) {Alg. I};
\node[right=6mm of a] (b) {Alg. II};
\node[right=6mm of b] (c) {Est. Datos I};
\node[right=6mm of c] (d) {Est. Datos II};
\node[right=6mm of d] (e) {Alg. y Compl.};
\node[below=6mm of e] (f) {Bases de Datos};
\node[left=6mm of f] (g) {Diseño SW I};
\node[left=6mm of g] (h) {Diseño SW II};
\node[left=6mm of h] (i) {Proyecto Final};
\foreach \p/\q in {a/b,b/c,c/d,d/e,e/f,f/g,g/h,h/i}{\draw[->,azul] (\p) -- (\q);}
\end{tikzpicture}}
\vfill
{\scriptsize\color{gris}Perder Algoritmia y Programación I bloquea 12 materias en cadena (BFS).}
\end{frame}

\begin{frame}{La regla de matrícula es una función booleana}
\unidad{Unidad 2 · álgebra booleana: expresión, Karnaugh y De Morgan}
\begin{columns}[T]
\begin{column}{0.58\textwidth}
\footnotesize
Con $A_x$: $x$ aprobada y $C_x$: $x$ cursando (ambas a la vez es imposible: \emph{don't care}):
\[ \mathrm{Disp}(C2)=\overline{A_{C2}}\cdot\overline{C_{C2}}\cdot A_{C1}\cdot(A_F+C_F) \]
\begin{itemize}
\item Forma canónica: $R=\Sigma m(5,6)+d(3,7)$
\item Karnaugh con \emph{don't cares}: $R=A_{C1}A_F+A_{C1}C_F$
\item De Morgan: $\overline R=\overline{A_{C1}}+\overline{A_F}\,\overline{C_F}$
\end{itemize}
\end{column}
\begin{column}{0.38\textwidth}
\centering\small Mapa de Karnaugh de $R$\\[2mm]
\begin{tabular}{c|cccc}
 & \multicolumn{4}{c}{$A_F\,C_F$}\\
$A_{C1}$ & 00 & 01 & 11 & 10\\\hline
0 & 0 & 0 & X & 0\\
1 & 0 & \cellcolor{azul!25}1 & \cellcolor{azul!25}X & \cellcolor{azul!25}1\\
\end{tabular}
\end{column}
\end{columns}
\vfill
{\scriptsize\color{gris}$\overline R$ lista lo que le falta al estudiante. Verificado en 62 materias y 17\,750 combinaciones: 0 diferencias con el grafo.}
\end{frame}

\begin{frame}{Importancia: las materias más críticas}
\unidad{Unidad 4 · centralidad: grado, alcance e intermediación}
\centering\small
\begin{tabular}{lrrrr}
\toprule
Materia & Créditos & Bloquea & Puente & Puntaje\\
\midrule
Estructura de Datos II & 3 & 6 & 0.014 & 0.67\\
Algoritmia y Programación II & 3 & 9 & 0.008 & 0.66\\
Estructura de Datos I & 3 & 7 & 0.011 & 0.64\\
Algoritmia y Programación I & 3 & 12 & 0.000 & 0.62\\
Cálculo I & 5 & 10 & 0.000 & 0.62\\
\bottomrule
\end{tabular}

\vspace{2mm}
$I(v)=0.2\,\hat c(v)+0.5\,\hat b(v)+0.3\,\hat p(v)$, cada componente dividido entre su máximo.
\vfill
{\scriptsize\color{gris}Bloqueo = alcance; puente = intermediación. Cercanía armónica: la red no es conexa. Pesos: decisión de modelación.}
\end{frame}

\begin{frame}{Validación: reproduce el promedio oficial}
\unidad{Caso real con resultado conocido}
\begin{columns}[c]
\begin{column}{0.4\textwidth}
\centering{\fontsize{60}{62}\selectfont\bfseries\color{acento}3.70}\\[2mm]
{\small PGA de la aplicación de la universidad (estudiante real, 27 créditos que cuentan)}
\end{column}
\begin{column}{0.56\textwidth}
\small
\begin{tabular}{lr}
\toprule
Cómo se cuenta & Resultado\\
\midrule
Todos los intentos, sin retiradas & \textbf{3.70}\\
Solo el último intento & 3.85\\
\bottomrule
\end{tabular}

\vspace{3mm}
72 pruebas automáticas y equivalencia lógica de la regla booleana.
\end{column}
\end{columns}
\vfill
{\footnotesize\color{gris}Las notas de la aplicación son una guía y no oficiales. Caso real de un integrante del grupo, sin nombre.}
\end{frame}

\begin{frame}{Demostración con el historial de validación}
\small
\begin{tabular}{p{0.4\textwidth}p{0.52\textwidth}}
\toprule
Pregunta & Respuesta\\
\midrule
¿Puedo matricular Cálculo III? & No: te falta aprobar Cálculo II (estado: retirada)\\[1mm]
¿Qué puedo matricular el próximo semestre? & 14 créditos sugeridos: Cálculo II, Física Calor-Ondas, POO y Estructura de Datos I\\[1mm]
¿Y si pierdo Algoritmia II? & El promedio baja de 3.70 a entre 3.48 y 3.62; sigue en estado normal\\[1mm]
¿Puedo retirarla? & Sí: quedarías con 10 créditos (el mínimo es 7)\\
\bottomrule
\end{tabular}
\vfill
{\scriptsize\color{gris}\texttt{python main.py --historial}\\\texttt{data/estudiantes/validacion\_pga.csv}}
\end{frame}

\begin{frame}{Perceptrón: método reproducido}
\unidad{Unidad 5 · sin datos reales de notas todavía}
\begin{columns}[T]
\begin{column}{0.5\textwidth}
\small\textbf{Regla de aprendizaje}\\[1mm]
$z=\sum_i x_i w_i+b,\quad \hat y=\varphi(z)$\\
$w\leftarrow w+\gamma\,(y-\hat y)\,x$\\
$b\leftarrow b+\gamma\,(y-\hat y)$\\
$\gamma=0.05$, $b_0=0.122$
\end{column}
\begin{column}{0.46\textwidth}
\small\textbf{AND de clase, 3 pasos}\\[1mm]
$b: 0.122\to 0.072$\\
$w_2=-0.05$\\
$w=(-0.05,-0.05),\ b=-0.028$
\end{column}
\end{columns}
\vspace{5mm}
\small Se aplica a predecir si una nota será $\ge 3.0$. No se recolectaron notas reales: queda como demostración del método y trabajo futuro.
\end{frame}

\begin{frame}{Conclusiones y limitaciones}
\begin{columns}[T]
\begin{column}{0.48\textwidth}
\textbf{Conclusiones}
\small\begin{itemize}
\item El plan es un DAG y las consultas se resuelven en tiempo lineal.
\item La regla booleana equivale exactamente al grafo.
\item El promedio coincide con el oficial en un caso real.
\item La centralidad señala las materias más críticas.
\end{itemize}
\end{column}
\begin{column}{0.48\textwidth}
\textbf{Limitaciones}
\small\begin{itemize}
\item Solo Ingeniería de Sistemas.
\item Encuesta por conveniencia ($n=43$).
\item Perceptrón sin datos reales de notas.
\item No modela la oferta real ni el límite por idiomas.
\end{itemize}
\end{column}
\end{columns}
\end{frame}

\begin{frame}[plain]
\vfill\centering{\fontsize{28}{30}\selectfont\bfseries Gracias. ¿Preguntas?}\\[8mm]
{\small Código, datos y README para reproducir el proyecto:}\\[2mm]
\url{https://github.com/USUARIO/REPOSITORIO} \pendiente{enlace real}
\vfill
\end{frame}

\end{document}
